\subsection{Extractie van detectiekenmerken uit e-mails}
\label{subsec:extractie-van-detectiekenmerken}

Zoals reeds besproken in sectie~\ref{subsubsec:Afzender en e-mailheaders} volgt een e-mail de standaard uit RFC~5322 \autocite{Resnick2008}. Een e-mail bestaat dus uit headers en een body, die meerdere onderdelen kan bevatten, zoals platte tekst, HTML en bijlagen. Een mailclient vertaalt dit onderliggende bericht naar een weergave die de gebruiker kan lezen en interpreteren. Echter kan daarbij informatie verloren gaan. Zo toont een mailclient vaak enkel de weergavenaam van de afzender en niet het volledige e-mailadres. Ook de zichtbare tekst van een hyperlink hoeft niet overeen te komen met de URL waarnaar de link daadwerkelijk verwijst \autocite{Veit2025}. De gebruiker ziet daardoor niet de werkelijke bestemming van de link, al kan deze wel uit de onderliggende HTML-code worden afgeleid. Volgens \textcite{Chen2020} kunnen aanvallers deze verschillen misbruiken om toch een legitiem ogende mail te fabriceren.

Een toepassing die een e-mail wil beoordelen, heeft dus het onderliggende bericht nodig en niet enkel de weergave ervan. Hieronder worden de mogelijke routes besproken:

\begin{enumerate}
    \item \textbf{\texttt{.eml}-bestanden.} Een \texttt{.eml}-bestand bevat een volledige e-mail in het formaat volgens RFC~5322 \parencite{Resnick2008}, inclusief de headers die de ontvangende server heeft toegevoegd. Dit formaat komt overeen met dat van de datasets uit sectie~\ref{subsubsec:datasets}.
    \item \textbf{Gmail API.} Mits toestemming van de gebruiker kan een toepassing een bericht rechtstreeks opvragen bij de mailprovider \footfullcite{GmailMessagesGet}. De e-mail kan zowel als origineel bericht of als een door Gmail gestructureerde weergave worden opgevraagd.
\end{enumerate}

De verkregen gegevens zijn niet altijd onmiddellijk bruikbaar voor kenmerkextractie, omdat de headers en de berichtinhoud verschillende tekensets en coderingen kunnen gebruiken \autocite{Koide2025}. Zo beschrijft RFC~2047 \autocite{Moore1996} hoe niet-ASCII-tekst in headers wordt weergegeven met behulp van zogenaamde \emph{encoded-words}, die moeten worden gedecodeerd om de oorspronkelijke tekst te achterhalen. Ook de wijze van aanlevering kan een codering toevoegen: de Gmail API levert het volledige bericht bij \texttt{format=raw} bijvoorbeeld als base64url-gecodeerde tekenreeks \footfullcite{GmailMessages}. Waar nodig worden de gegevens daarom eerst gedecodeerd voordat kenmerken eruit kunnen worden afgeleid.